% Physical unit 035; accepted Session 03 after fresh 400-dpi authority replay.
\GordanUnitStart{035}{printed page 30}
\GordanTopFolio{30}
\begingroup\small
\[
\begin{aligned}
&={1\over\sin A}\cdot{1-2q\over1+q^2}\cdot
 {1-3q^2+4q^2\cos^2{iv\pi\over2K}
  \over 1+2q-4q\cos^2{iv\pi\over2K}}\\
&={1\over\sin A}\cdot{1-2q\over1+q^2}\cdot
 {1-3q^2+q\varepsilon\over 1+2q-\varepsilon{1-4q^2\over1-2q\varepsilon}}\\
&={\sqrt{x'}\over\sin A}\cdot{1-2q\varepsilon\over1+q^2}\cdot {1+q\varepsilon-3q^2\over1-\varepsilon}.
\end{aligned}
\]
Or, since
\[
 Civ={1\over C(v,x')}={1\over \sin l_0},
\]
there results
\[
 \sin A={\sqrt{x'}\sin l_0\over1-\varepsilon}\{1-q\varepsilon-4q^2\}
 =\sqrt{\sin l_0\sin B_0}\left\{1+{\varepsilon^2\over2}-q\varepsilon-4q^2\right\}.
\]

\sect{9}

Now, however, it is useful to establish the other formula for computing the longitude of the period. For there is:
\[
\begin{aligned}
1-{2\over\pi}\mathcal A(l_0)&={2K\over\pi}iZ(iv+K)\qquad \text{cap. VI \S\ 4.}\\
&={-4iq\sin {iv\pi\over K}\over 1-2q\cos {iv\pi\over K}}
=-{8iq\tg {iv\pi\over2K}\cos^2{iv\pi\over2K}
\over 1+2q-4q\cos^2{iv\pi\over2K}}\\
&={8q\cos A\over1+2q\sin^2A-2q}
={2\varepsilon\cos A\over {1-2q\varepsilon\over1-2q}-\varepsilon}
=2\varepsilon{1-2q\over1-\varepsilon}\cos A.
\end{aligned}
\]
\[
 \mathcal A(l_0)={\pi\over2}-\pi\varepsilon{1-2q\over1-\varepsilon}\cos A.
\]

\rule{0.35\textwidth}{0.4pt}

\caput{VIII}
\sect{1}

In the sixth chapter we showed that the quantities
\[
 B,\quad l,\quad T,\quad s,\quad \psi,\quad \varphi,
\]
can be determined by functions of one variable $u$. In order to bring them to calculation, the $\Theta$ functions are to be developed into series, proceeding according to the ratio and powers of the constant $q$ and the cosines and sines of multiples of the fraction
\[
 {\pi u\over2K},
\]
which, in the manner of analysts, we denote by \gq{$x$}.

% Continuation: source PDF pages 036--047. English translation.
\endgroup
